\documentclass[10pt,a4paper]{amsart}
\usepackage[margin=19mm]{geometry}
\usepackage{amsmath,amssymb,mathtools}
\usepackage[scr=rsfs,cal=euler]{mathalfa}
\usepackage{xfrac}
\usepackage[unicode,hidelinks,bookmarks=false]{hyperref}
\newcommand{\ie}{i.e.\ }
\newcommand{\cf}{cf.\ }
\newcommand{\resp}{respectively\ }
\newcommand{\Subsection}{\subsection}
\providecommand{\nobreakdash}{\nobreak\hbox{-}}
\setlength{\emergencystretch}{2em}
\multlinegap0pt
\usepackage{bm}
\usepackage{bbm}
\usepackage{dsfont}
\usepackage{xspace}
\usepackage{cancel}
\usepackage{mathtools}
\usepackage{cleveref}
\usepackage{amsthm}
\makeatletter
\@ifundefined{theorem}{\newtheorem{theorem}{Theorem}}{}
\@ifundefined{assumption}{\newtheorem{assumption}[theorem]{Assumption}}{}
\@ifundefined{lemma}{\newtheorem{lemma}[theorem]{Lemma}}{}
\makeatother
\title[Coordinate-energy transformation and the one-point function ]{Coordinate-energy transformation and the one-point function for the Heisenberg-Ising XXZ spin-1/2 chain on the ring}
\author{Eric I. Corwin, Nikolaus Elsaesser,, Axel Saenz}
\date{Source: arXiv:2506.14171v3 (CC BY 4.0)}
\begin{document}
\maketitle

\noindent\emph{Source and licence.} Coordinate-energy transformation and the one-point function for the Heisenberg-Ising XXZ spin-1/2 chain on the ring. Version arXiv:2506.14171v3 (CC BY 4.0), distributed under the Creative Commons Attribution 4.0 International licence, as recorded in the arXiv licence metadata for this exact version and shown on the abstract page https://arxiv.org/abs/2506.14171v3. Full rights notice: licenses/arxiv-2506.14171-CC-BY-4.0.txt.\par\smallskip

\noindent\textit{Editorial scope.} Counted unit: the Lemma whose source label is \texttt{l:Indicator function expansion} in the article (source0.tex lines 1256--1262, source label \texttt{l:Indicator function expansion}), delivered together with the article's own complete original proof of it (source0.tex lines 1265--1532, one \texttt{proof} environment of 268 lines). No mathematics is written, repaired or re-derived: statement and proof are the article's own text line for line. Required context retained uncounted: a:N=2 proof assumptions (lines 1200--1202); e:contour (lines 1224--1228); e:integrand (lines 1230--1235); l: xi\_1 poles simple (lines 1726--1728). Every definition, display and label of those lines is retained unchanged. Omitted from this package and not redistributed: 1--1199, 1203--1223, 1229--1229, 1236--1255, 1263--1264, 1533--1725, 1729--1980. Nothing inside a retained excerpt is omitted or altered: every retained body is delivered exactly as the article's lines are written, so no omission receipt of any kind accompanies this package. Citations: the retained excerpts carry no citation command at all, so no bibliography entry of the article is transcribed and no reference is redistributed. Reference adaptation: none was needed and none was made; every reference inside the delivered unit resolves inside it, so no unresolved reference or citation is produced. Printed slips kept literal: no printed word, symbol, hypothesis or number of the counted statement or of its proof was altered. Layout and class: the article's own class is \texttt{article} with packages including inputenc, textcomp, newunicodechar, geometry, amsmath, amssymb, amsthm, mathtools, graphicx, hyperref, caption, cleveref, subcaption, algorithm; the wrapper is the lane's pinned \texttt{amsart} editorial wrapper at 10pt on A4, which restates the theorem-like environments the retained text needs and adds the article's own macro definitions that the retained text needs (none), with extra wrapper packages (bm, bbm, dsfont, xspace, cancel, mathtools, cleveref). The delivered numbering is the unit's own. Assets: no retained span contains a figure environment, a graphics-inclusion command, a PDF-page inclusion, a table or a TikZ picture; no image file is copied into this package, the package declares no asset dependency and no image file is redistributed with it. Licence: version arXiv:2506.14171v3 of this article is distributed under the Creative Commons Attribution 4.0 International licence, as recorded in the arXiv licence metadata for this exact version and shown on the abstract page https://arxiv.org/abs/2506.14171v3. A full rights notice is delivered at licenses/arxiv-2506.14171-CC-BY-4.0.txt. Characters: every retained excerpt is pure ASCII, so the delivered engine, XeLaTeX with its default Latin Modern text font, reports no missing character.
\begin{assumption}\label{a:N=2 proof assumptions}
    Assume $N=2$, $L>2$, $|\Delta|<\frac{L-1}{2L}$, $\Delta\neq\cos(2\pi k/(L-2))$ with $k\in [L-2]$, and that $\Delta$ is not in a finite set of values as determined in the proof of \Cref{l: xi_1 poles simple}.
\end{assumption}

\begin{equation}\label{e:contour}
\begin{split}
    C_j &= C_{r_j} \sqcup C_{R_j} \\
\end{split}
\end{equation}

\begin{equation}\label{e:integrand}
\begin{split}
    I(\boldsymbol{\xi};\mathbf{x},\mathbf{y})&= \frac{u(\xi_1,\xi_2;x_1,x_2)\xi_1^{-y_1-1}\xi_2^{-y_2-1}}{\left(\xi_1^L - S(\xi_1,\xi_2) \right)\left(\xi_2^L - S(\xi_2,\xi_1) \right)} \\
    &=  \frac{\xi_1^{x_1-y_1-1}\xi_2^{x_2-y_2-1}+S(\xi_2,\xi_1)\xi_1^{x_2-y_1-1}\xi_2^{x_1-y_2-1}}{\left(\xi_1^L - S(\xi_1,\xi_2) \right)\left(\xi_2^L - S(\xi_2,\xi_1) \right)}
\end{split}
\end{equation}

\begin{lemma}\label{l: xi_1 poles simple}
    The poles from $\xi_1^L-S(\xi_1,z_2)=0$ are simple except for finitely many values of $\Delta$. 
\end{lemma}

\begin{lemma}\label{l:Indicator function expansion}
Let $\boldsymbol{x},\boldsymbol{y} \in \mathcal{X}'$ and take \Cref{a:N=2 proof assumptions} to be true. Then, we have the following contour integral representation on the indicator function of equivalent configurations,
\begin{equation}\label{e:N=2ProofIndicator}
    \mathds{1}(\boldsymbol{x}=_{L} \boldsymbol{y}) = \oint_{C_1} d\xi_1 \oint_{C_2} d\xi_2 \, I(\boldsymbol{\xi};\mathbf{x},\mathbf{y}),
\end{equation}
where the contours $C_j$ are given by \eqref{e:contour} and the integrand $I$ is given by \eqref{e:integrand}.
\end{lemma}

\begin{proof}
We denote


\begin{equation}
\begin{split}
    &\mathds{1}(x=_L y) \coloneq \sum_{k\in\mathbb{Z}}\mathds{1}\left(x_1+kL=y_1,x_2+kL=y_2 \right) + \mathds{1}\left(x_2+(k-1)L=y_1,x_1+kL=y_2 \right).
\end{split}
\end{equation}
We claim that the non-permuted terms in the integrand will yield the indicator functions where both $x_1$ and $x_2$ are shifted by $kL$, and we claim that the permuted terms will yield the indicator functions where just $x_1$ and just $x_2$ are shifted by a factor of $L$ apart. That is,   
\begin{equation}\label{e:N=2ProofIdentityPart}
    \oint_{C_1}d\xi_1\oint_{C_2}d\xi_2 \frac{\xi_1^{x_1-y_1-1}\xi_2^{x_2-y_2-1}}{\left(\xi_1^L - S(\xi_1,\xi_2) \right)\left(\xi_2^L - S(\xi_2,\xi_1) \right)} = \sum_{k\in\mathbb{Z}}\mathds{1}(\mathbf{x}+kL=\mathbf{y})
\end{equation}
and 
\begin{equation}\begin{split}
    \oint_{C_1}d\xi_1\oint_{C_2}d\xi_2 \frac{S(\xi_2,\xi_1)\xi_1^{x_2-y_1-1}\xi_2^{x_1-y_2-1}}{\left(\xi_1^L - S(\xi_1,\xi_2) \right)\left(\xi_2^L - S(\xi_2,\xi_1) \right)} = \sum_{k\in\mathbb{Z}}\mathds{1}\left(x_2+(k-1)L=y_1,x_1+kL=y_2 \right).
    \end{split}
\end{equation}
Moreover, integrating with respect to the $C_{r_j}$ and $C_{R_j}$ contours separately for $\xi_1$ and $\xi_2$ will yield the terms of the sum over non-negative and negative $k$ terms respectively. 

Let's begin with the nested integrals over the $C_{r_j}$ contours and non-permuted terms. Because $r(C_{r_1})=r_j<1$, we have that $\left|\xi_i^L S(\xi_j,\xi_i) \right| \approx (r_j)^L<1$, so we can rewrite the denominator in the left hand side of \eqref{e:N=2ProofIdentityPart} as geometric sums\footnote{Note: the property $S(\xi_1,\xi_2)=(S(\xi_2,\xi_1))^{-1}$ will be frequently used without specific mention in this proof sketch.}. That is, 
\begin{equation}\label{e:N=2proofGeometricSum}
\begin{split}
    &\oint_{C_{r_1}}d\xi_1\oint_{C_{r_2}}d\xi_2 \frac{\xi_1^{x_1-y_1-1}\xi_2^{x_2-y_2-1}}{\left(\xi_1^L - S(\xi_1,\xi_2) \right)\left(\xi_2^L - S(\xi_2,\xi_1) \right)}\\ = &\sum_{k,m=0}^\infty \oint d\xi_1\oint d\xi_2\hspace{0.1cm}\xi_1^{x_1-y_1-1+kL}\xi_2^{x_2-y_2-1+mL}S(\xi_2,\xi_1)^{k-m}.
\end{split}
\end{equation}

To handle the sum, we will consider three cases: when $k>m$, when $k<m$, and $k=m$. For the first two cases, we will use the two possible change of variables for $\eta=\xi_1\xi_2$. Given $k<m$, we let $\xi_1=\eta/\xi_2$. Then we have that the corresponding terms from $\eqref{e:N=2proofGeometricSum}$ are equal to 
\begin{equation}
    \begin{split}
        \sum_{m>k}^\infty\oint_{C_{r_1\cdot r_2}}d\eta\oint_{C_{r_2}}d\xi_2 \ \eta^{x_1-y_1-1+kL}\xi_2^{x_2-x_1+y_1-y_2-1+(m-k)(L-1)}\left( \frac{\xi_2+\eta\xi_2-2\Delta\eta}{1+\eta-2\Delta\xi_2} \right)^{m-k}.
    \end{split}
\end{equation}
Focusing on $\xi_2$, our only possible pole inside the $C_{r_2}$ contour is at $\xi_2=0$ if its exponent is negative. Inspecting the exponent, we find that 
\begin{equation}
    x_2-x_1+y_2-y_1-1+(m-k)(L-1)\geq -(L-1)+(m-k)(L-1)\geq 0
\end{equation}
since $L>1$ and $m>k$. Thus there are no poles with respect to the the $\xi_2$ variable and the integrals in the $k<m$ case are zero. For $k>m$, we let $\xi_2=\eta/\xi_1$ and have a similar story. The corresponding terms in this case from \eqref{e:N=2proofGeometricSum} are equal to 
\begin{equation}
    \begin{split}
        \sum_{k>m}^\infty\oint_{C_{r_1\cdot r_2}}d\eta\oint_{C_{r_2}}d\xi_2 \ \xi_1^{x_1-x_2+y_2-y_1-1+(k-m)(L-1)}\eta^{x_2-y_2-1+mL}\left( \frac{\xi_1+\eta\xi_1-2\Delta\eta}{1+\eta-2\Delta\xi_1} \right)^{k-m}.
    \end{split}
\end{equation}

With a possible pole at $\xi_1=0$, we inspect the exponent of $\xi_1$ and find that 
\begin{equation}
    x_1-x_2+y_2-y_1-1+(k-m)(L-1)\geq -(L-1)+(k-m)(L-1)\geq0.
\end{equation}
Thus there are no poles with respect to the $\xi_1$ variable and the integrals in the $k>m$ case are zero. Now from \eqref{e:N=2proofGeometricSum} we are left with 
\begin{equation}
    \sum_{k=0}^{\infty}\oint_{C_{r_1}} d\xi_1\oint_{C_{r_2}} d\xi_2 \ \xi_1^{x_1-y_1-1+kL}\xi_2^{x_2-y_2-1+kL}.
\end{equation}
For a fixed $k\in\mathbb{Z}_{\geq 0}$, if $x_j+kL>y_j$ for $j=1$ or $j=2$, then there are no poles in the above integrand in $\xi_1$ or $\xi_2$ respectively, and the integral is $0$. If $x_j+kL = y_j$ for $j=1$ and $j=2$, have simple poles in $\xi_1$ and $\xi_2$ and a residue computation shows that
\begin{equation}
    \oint_{C_{r_1}} d\xi_1\oint_{C_{r_2}} d\xi_2 \ \xi_1^{x_1-y_1-1+kL}\xi_2^{x_2-y_2-1+kL} = \oint d\xi_1\oint d\xi_2 \ \xi_1^{-1}\xi_2^{-1} =1.
\end{equation}
If $x_j+kL<y_j$ for $j=1$ or $j=2$, there is a pole of order $-(x_j-y_j-1+kL)$ in $\xi_j$ and a residue computation shows that
\begin{equation}
     \oint_{C_{r_1}} d\xi_1\oint_{C_{r_2}} d\xi_2 \ \xi_1^{x_1-y_1-1+kL}\xi_2^{x_2-y_2-1+kL} = 0.
\end{equation}
Thus, we have that 
\begin{equation}
    \oint_{C_{r_1}}d\xi_1\oint_{C_{r_2}}d\xi_2 \frac{\xi_1^{x_1-y_1-1}\xi_2^{x_2-y_2-1}}{\left(\xi_1^L - S(\xi_1,\xi_2) \right)\left(\xi_2^L - S(\xi_2,\xi_1) \right)} = \sum_{k=0}^\infty \mathds{1}(x+kL = y)
\end{equation}
where $x+kL = (x_1+kL, x_2+kL)$.

Turning to the permuted terms, since $|\xi_j|< 1$ in the $C_{r_j}$ contours, we have
\begin{equation}\begin{split}
&\oint_{C_{r_1}}d\xi_1\oint_{C_{r_2}}d\xi_2 \frac{S(\xi_2,\xi_1)\xi_1^{x_2-y_1-1}\xi_2^{x_1-y_2-1}}{\left(\xi_1^L - S(\xi_1,\xi_2) \right)\left(\xi_2^L - S(\xi_2,\xi_1) \right)}\\
&=\sum_{k,m=0}^\infty \oint_{C_{r_1}}d\xi_1\oint_{C_{r_2}}d\xi_2 \xi_1^{x_2-y_1-1}\xi_2^{x_1-y_2-1}S(\xi_2,\xi_1)^{k-m+1}.
\end{split}
\end{equation}
Now we will split the sum into three possible cases: $k-m+1<0$, $k-m+1>0$, and $k-m+1=0$. For $k-m+1<0$, we use the change of variables $\xi_1=\frac{\eta}{\xi_2}$ which gives us
\begin{equation}
    \begin{split}
        &\sum_{k-m+1<0}\oint_{C_{r_1\cdot r_2}}d\eta\oint_{C_{r_2}}d\xi_2 \ \eta^{x_2-y_1-1+kL}\xi_2^{x_1-x_2+y_1-y_2-1+(m-k)L}\left(\frac{1+\eta-2\Delta\eta/\xi_2}{1+\eta-2\Delta\xi_2} \right)^{m-k-1} \\
        &= \sum_{k-m+1<0}\oint_{C_{r_1\cdot r_2}}d\eta\oint_{C_{r_2}}d\xi_2 \ \eta^{x_2-y_1-1+kL}\xi_2^{x_1-x_2+y_1-y_2+(m-k)L-(m-k)}\left(\frac{\xi_2+\eta\xi_2-2\Delta\eta}{1+\eta-2\Delta\xi_2} \right)^{m-k-1}.
    \end{split}
\end{equation}
Since $m-k>1$ in this case, we have the following lower bound on the $\xi_2$ exponent $x_1-x_2+y_1-y_2+(m-k)(L-1)\geq -2(L-1)+(m-k)(L-1)\geq 0$. Thus, there are no poles in the $\xi_2$ integral and the terms in this case are all zero. For $k-m+1>0$, we use the change of variables $\xi_2=\frac{\eta}{\xi_1}$. Then we have 
\begin{equation}
    \begin{split}
        &\sum_{k-m+1>0}\oint_{C_{r_1}}d\xi_1\oint_{C_2}d\eta \ \xi_1^{x_2-x_1+y_2-y_1-1+(k-m)L}\eta^{x_1--y_2-1+mL}\left(\frac{1+\eta-2\Delta\eta/\xi_1}{1+\eta-2\Delta\xi_1} \right)^{k-m+1} \\
        &=\sum_{k-m+1>0}\oint_{C_{r_1}}d\xi_1\oint_{C_{r_2}}d\eta \ \xi_1^{x_2-x_1+y_2-y_1+(k-m)L-(k-m)-2}\eta^{x_1--y_2-1+mL}\left(\frac{\xi_1+\eta\xi_1-2\Delta\eta}{1+\eta-2\Delta\xi_1} \right)^{k-m+1}.
        \end{split}
\end{equation}
Since $k-m\geq0$ in this case, we have the following lower bound on the $\xi_1$ exponent $x_2-x_1+y_2-y_1-2+(k-m)L-(k-m)\geq (k-m)(L-1) \geq 0$. Hence there are no poles in the $\xi_1$ integral, and the terms in this case are zero. Finally for $k-m+1 = 0$ we have the remaining terms
\begin{equation}
    \sum_{k=0}^\infty\oint_{C_{r_1}}d\xi_1\oint_{C_{r_2}}d\xi_2\ \xi_1^{x_2-y_1-1+kL}\xi_2^{x_1-y_2-1+(k+1)L}.
\end{equation}
If $x_1 + (k+1)L=y_2$ and $x_2+kL=y_1$, then we have 
\begin{equation}
    \oint_{C_{r_1}}d\xi_1\oint_{C_{r_2}}d\xi_2 \ \xi_1^{-1}\xi_2^{-1} = 1.
\end{equation}
If $x_1 + (k+1)L<y_2$ or $x_2+kL<y_1$, then a residue computation yields zero. And if $x_1 + (k+1)L>y_2$ or $x_2+kL>y_1$, then $\xi_1$ or $\xi_2$ do not have a pole and the integral is zero. Thus, 
\begin{equation}
\begin{split}
    \oint_{C_{r_1}}d\xi_1\oint_{C_{r_2}}d\xi_2 \frac{S(\xi_2,\xi_1)\xi_1^{x_2-y_1-1}\xi_2^{x_1-y_2-1}}{\left(\xi_1^L - S(\xi_1,\xi_2) \right)\left(\xi_2^L - S(\xi_2,\xi_1) \right)}&=\sum_{k=0}^\infty\mathds{1}\left(x_2+kL=y_1,\ x_1+(k+1)L=y_2 \right)\\
    &=\sum_{k=1}^\infty\mathds{1}\left(x_2+(k-1)L=y_1,\ x_1+kL=y_2 \right)
    \end{split}
\end{equation}

Our integration over the $C_{R_j}$ contours of $I(\xi;x,y)$ will be nearly identical to that of the $C_{r_j}$ contours. The key difference is that we obtain geometric sums using the fact that $|\xi_j^{-L} S(\xi_j,\xi_i)|\approx (r_j)^L <1$ on $C_{R_j}$ and then we can swap the sum over non-negative integers to one over negative integers, and then the integrands are essentially the same as in the case of integrating over $C_{r_j}$ contours. Beginning with non permuted terms, we have
\begin{equation}
    \begin{split}
        &\oint_{C_{R_1}}d\xi_1\oint_{C_{R_2}}d\xi_2 \frac{\xi_1^{x_1-y_1-1}\xi_2^{x_2-y_2-1}}{\left(\xi_1^L - S(\xi_1,\xi_2) \right)\left(\xi_2^L - S(\xi_2,\xi_1) \right)} \\ 
        &=\sum_{k,m=0}^\infty \oint_{C_{R_1}}d\xi_1\oint_{C_{R_2}}d\xi_2 \ \xi_1^{x_1-y_1-1-(k+1)L}\xi_2^{x_2-y_2-1-(m+1)L}S(\xi_2,\xi_1)^{m-k} \\ 
        &= \sum_{k,m\leq0} \oint_{C_{R_1}}d\xi_1\oint_{C_{R_2}}d\xi_2 \ \xi_1^{x_1-y_1-1+(k-1)L}\xi_2^{x_2-y_2-1+(m-1)L}S(\xi_2,\xi_1)^{k-m}.
    \end{split}
\end{equation}

As we did previously, we consider three cases: $k<m$, $k>m$, and $k=m$. For $k<m$, we use the change of variables $\xi_2=\frac{\eta}{\xi_1}$. Then we have 
\begin{equation}
    \begin{split}
        &\sum_{k<m\leq 0} \oint_{C_{R_1}}d\xi_1\oint_{C_{R_1\cdot R_2}}d\eta \ \xi_1^{x_1-x_2+y_2-y_1-1+(k-m)L} \eta^{x_2-y_2-1+(m-1)L}\left(\frac{1+\eta-2\Delta\xi_1}{1+\eta-2\Delta\eta/\xi_1} \right)^{m-k} \\
        &=\sum_{k<m\leq0} \oint_{C_{R_1}}d\xi_1\oint_{C_{R_1\cdot R_2}}d\eta \ \xi_1^{x_1-x_2+y_2-y_1-1+(k-m)(L-1)} \eta^{x_2-y_2-1+(m-1)L}\left(\frac{1/\xi_1+\eta/\xi_1-2\Delta}{1+\eta-2\Delta\eta/\xi_1} \right)^{m-k}.
    \end{split}
\end{equation}
 Since $r(C_{R_1})>1$, potential poles for $\xi_1$ exist at $+\infty$ when its exponent is positive. Since $k-m<0$ in this case, we have the following upper bound on the exponent of $\xi_1$: $x_1-x_2+y_2-y_1-1+(k-m)(L-1)\leq -2 + L-1 +(k-m)(L-1) \leq -2$. Hence there are no poles with respect to $\xi_1$ and the integrals in this case are zero. For $k>m$ we let $\xi_1=\frac{\eta}{\xi_2}$. Then we have 
\begin{equation}
    \begin{split}
        &\sum_{0\geq k>m} \oint_{C_{R_1 \cdot R_2}}d\eta\oint_{C_{R_2}}d\xi_2 \ \eta^{x_1-y_1-1+kL} \xi_2^{x_2-x_1+y_1-y_2-1+(m-k)L}\left(\frac{1+\eta-2\Delta\xi_2}{1+\eta-2\Delta\eta/\xi_2} \right)^{k-m} \\
        &= \sum_{0\geq k>m} \oint_{C_{R_1\cdot R_1}}d\eta\oint_{C_{R_2}}d\xi_2 \ \eta^{x_1-y_1-1+kL} \xi_2^{x_2-x_1+y_1-y_2-1+(m-k)(L-1)}\left(\frac{1/\xi_2+\eta/\xi_2-2\Delta}{1+\eta-2\Delta\eta/\xi_2} \right)^{k-m}.
    \end{split}
\end{equation}
Since $m-k<0$, we have the following upper bound on the $\xi_2$ exponent: $x_2-x_1+y_1-y_2-1+(m-k)(L-1)\leq L-1 - 2 + (m-k)(L-1) \leq -2$. Thus there are no poles with respect to $\xi_2$ and the integrals in this case are zero. Lastly we have the case where $k=m$. Then we have the remaining terms 
\begin{equation}
    \begin{split}
        \sum_{k\leq 0}\oint_{C_{R_1}}d\xi_1 \oint_{C_{R_2}}d\xi_2 \ \xi_1^{x_1-y_1 - 1 +(k-1)L}\xi_2^{x_2-y_2-1+(k-1)L}.
    \end{split}
\end{equation}

If $x_j + (k-1)L = y_j$ for $j=1,2$, we have 
\begin{equation}
    \oint_{C_{R_1}}d\xi_1\oint_{C_{R_2}}d\xi_2  \ \xi_1^{-1}\xi_2^{-1} = 1.
\end{equation}
If $x_j - y_j + (k-1)L>0$ for $j=1$ or $j=2$ then a quick residue computation yields zero for the integrals in this case. And if $x_j-y_j+(k-1)L<0$ for $j=1$ or $j=2$, then there are no poles for $\xi_1$ or $\xi_2$, and the integrals in this case are also zero. Thus, 
\begin{equation}
    \oint_{C_{R_1}}d\xi_1\oint_{C_{R_2}}d\xi_2 \frac{\xi_1^{x_1-y_1-1}\xi_2^{x_2-y_2-1}}{\left(\xi_1^L - S(\xi_1,\xi_2) \right)\left(\xi_2^L - S(\xi_2,\xi_1) \right)}=\sum_{k\leq -1}\mathds{1}\left(x_1+kL=y_1, \, x_2+kL=y_2\right).
\end{equation}

Following the same process for the permuted terms, we begin with
\begin{equation}
    \begin{split}
        &\oint_{C_{R_1}}d\xi_1\oint_{C_{R_2}}d\xi_2 \frac{S(\xi_2,\xi_1)\xi_1^{x_2-y_1-1}\xi_2^{x_1-y_2-1}}{\left(\xi_1^L - S(\xi_1,\xi_2) \right)\left(\xi_2^L - S(\xi_2,\xi_1) \right)} \\
        &= \sum_{k,m\leq 0} \oint_{C_{R_1}}d\xi_1\oint_{C_{R_2}}d\xi_2 \xi_1^{x_2-y_1-1+(k-1)L}\xi_2^{x_1-y_2-1+(m-1)L}S(\xi_2,\xi_1)^{k-m+1}.
    \end{split}
\end{equation}

We separate the sum into three cases: $k-m+1>0$, $k-m+1<0$, and $k-m+1 = 0$. For $k-m+1>0$ we use the change of variables $\xi_1=\frac{\eta}{\xi_2}$ and have 
\begin{equation}
    \begin{split}
        &\sum_{0\geq k>m-1} \oint_{C_{R_1 \cdot R_2}}d\eta\oint_{C_{R_2}}d\xi_2 \ \eta^{x_2-y_1-1+(k-1)L}\xi_2^{x_1-x_2+y_1-y_2-1+(m-k)L}\left(\frac{1+\eta-2\Delta\xi_2}{1+\eta-2\Delta\eta/\xi_2}\right)^{k-m+1} \\
        &=\sum_{0\geq k>m-1}\oint_{C_{R_1\cdot R_2}}d\eta\oint_{C_{R_2}}d\xi_2 \ \eta^{x_2-y_1-1+(k-1)L}\xi_2^{x_1-x_2+y_1-y_2+(m-k)(L-1)}\left(\frac{1/\xi_2+\eta/\xi_2-2\Delta}{1+\eta-2\Delta\eta/\xi_2}\right)^{k-m+1}.
    \end{split}
\end{equation}

Since $m-k<1$, we have the following upper bound on the $\xi_2$ exponent: $x_1-x_2 +y_1-y_2+(m-k)(L-1)\leq -2 + (m-k)(L-1) \leq -2$. Thus there are no poles with respect to $\xi_2$ and the integrals in this case are zero. For $k-m+1<0$ we use the change of variables $\xi_2=\frac{\eta}{\xi_1}$ and have 
\begin{equation}
    \begin{split}
        &\sum_{0\geq m>k+1}\oint_{C_{R_1}}d\xi_1\oint_{C_{R_1\cdot R_2}}d\eta \ \xi_1^{x_2-x_1+y_2-y_1-1+(k-m)L}\eta^{x_1-y_2-1+(m-1)L}\left(\frac{1+\eta-2\Delta\xi_1}{1+\eta-2\Delta\eta/\xi_1}\right)^{m-k-1} \\
        &=\sum_{0\geq m>k+1}\oint_{C_{R_1}}d\xi_1\oint_{C_{R_1\cdot R_2}}d\eta \ \xi_1^{x_2-x_1+y_2-y_1-2+(k-m)(L-1)}\eta^{x_1-y_2-1+(m-1)L}\left(\frac{1/\xi_1+\eta/\xi_1-2\Delta}{1+\eta-2\Delta\eta/\xi_1}\right)^{m-k-1}.
    \end{split}
\end{equation}
Since $k-m<-1$, we have the following upper bound on the $\xi_1$ exponent: $x_2-x_1+y_2-y_1-2 + (k-m)(L-1) \leq 2(L-1) + (k-m)(L-1)-2 \leq -2$. Hence there are no poles with respect to $\xi_1$ and the integrals in this case are zero. For $k-m+1=0$, we have
\begin{equation}
    \sum_{k\leq -1}\oint_{C_{R_1}}d\xi_1\oint_{C_{R_2}}d\xi_2 \ \xi_1^{x_2-y_1-1+(k-1)L}\xi_2^{x_1-y_2-1+kL}.
\end{equation}

If $x_2+(k-1)L=y_1$ and $x_1+kL=y_2$, then we have 
\begin{equation}
    \oint_{C_{R_1}}d\xi_1\oint_{C_{R_2}}d\xi_2 \ \xi_1^{-1}\xi_2^{-1} = 1.
\end{equation}

Otherwise, we have negative exponents that yield no poles with respect to $\xi_1$ or $\xi_2$, or we have positive exponents on both $\xi_1$ and $\xi_2$ and straightforward residue computations yield zero. Thus, 
\begin{equation}
    \oint_{C_{R_1}}d\xi_1\oint_{C_{R_2}}d\xi_2 \frac{S(\xi_2,\xi_1)\xi_1^{x_2-y_1-1}\xi_2^{x_1-y_2-1}}{\left(\xi_1^L - S(\xi_1,\xi_2) \right)\left(\xi_2^L - S(\xi_2,\xi_1) \right)} = \sum_{k\leq -1}\mathds{1}\left(x_2+(k-1)L=y_1, \ x_1+kL=y_2 \right).
\end{equation}

Up to this point, we have shown 
\begin{equation}
    \begin{split}
        &\oint_{C_{r_1}}d\xi_1\oint_{C_{r_2}}d\xi_2 \ \frac{\xi_1^{x_1-y_1-1}\xi_2^{x_2-y_2-1}+ S(\xi_2,\xi_1)\xi_1^{x_2-y_1-1}\xi_2^{x_1-y_2-1}}{\left( \xi_1^L-S(\xi_1,\xi_2) \right)\left( \xi_2^L-S(\xi_2,\xi_1) \right)} \\
        &= \sum_{k=0}^\infty \mathds{1}\left( x_1+kL=y_1,  \ x_2+kL=y_2 \right) + \sum_{k=1}^\infty\mathds{1}\left( x_2+(k-1)L=y_1,  \ x_1+kL=y_2 \right) \\ \text{and}\\
        &\oint_{C_{R_1}}d\xi_1\oint_{C_{R_2}}d\xi_2 \ \frac{\xi_1^{x_1-y_1-1}\xi_2^{x_2-y_2-1}+ S(\xi_2,\xi_1)\xi_1^{x_2-y_1-1}\xi_2^{x_1-y_2-1}}{\left( \xi_1^L-S(\xi_1,\xi_2) \right)\left( \xi_2^L-S(\xi_2,\xi_1) \right)}\\
        &= \sum_{k\leq -1} \mathds{1}\left( x_1+kL=y_1,  \ x_2+kL=y_2 \right) + \mathds{1}\left( x_2+(k-1)L=y_1,  \ x_1+kL=y_2 \right).
    \end{split}
\end{equation}

Integrating over the $C_{r_1}$ contour for $\xi_1$ and $C_{R_2}$ contour for $\xi_2$, we begin with
\begin{equation}
    \begin{split}
        &\oint_{C_{r_1}}d\xi_1\oint_{C_{R_2}}d\xi_2 \ \frac{\xi_1^{x_1-y_1-1}\xi_2^{x_2-y_2-1}}{\left( \xi_1^L-S(\xi_1,\xi_2) \right)\left( \xi_2^L-S(\xi_2,\xi_1) \right)} \\
        &=-\sum_{k,m=0}^{\infty}\oint_{C_{r_1}}d\xi_1\oint_{C_{R_2}}d\xi_2 \ \xi_1^{x_1-y_1-1+kL}\xi_2^{x_2-y_2-1-(m+1)L}S(\xi_2,\xi_1)^{k+m+1} 
    \end{split}
\end{equation}
since $|\xi_1^LS(\xi_2,\xi_1)|\approx \frac23|\Delta|\epsilon^{L-1}< 1$ and $|\xi_2^{-L}S(\xi_2,\xi_1)|\approx \frac23|\Delta|(2\epsilon)^{L-1}< 1$. For $\xi_2$, we have potential poles at $+\infty$. Using the change of variables $\xi_1=\frac{\eta}{\xi_2}$, we get 
\begin{equation}
    \sum_{k,m=0}^{\infty}\oint_{C_{r=1/2}}d\eta\oint_{C_{R_2}}d\xi_2 \ \eta^{x_1-y_1-1+kL}\xi_2^{x_2-x_1+y_1-y_2-(k+m+1)(L-1)-1}\left(\frac{1/\xi_2+\eta/\xi_2-2\Delta}{1+\eta-2\Delta\eta/\xi_2}\right)^{k+m+1}.
\end{equation}
Then we have the following upper bound on the $\xi_2$ exponent: $x_2-x_1+y_1-y_2-(k+m+1)(L-1)-1\leq L-1-2-(k+m+1)(L-1) \leq -2$. Hence there are no poles with respect to $\xi_2$ and the above integrals are zero. Moving along, we compute the permuted terms
\begin{equation}
    \begin{split}
        &\oint_{C_{r_1}}d\xi_1\oint_{C_{R_2}}d\xi_2 \ \frac{S(\xi_2,\xi_1)\xi_1^{x_2-y_1-1}\xi_2^{x_1-y_2-1}}{\left( \xi_1^L-S(\xi_1,\xi_2) \right)\left( \xi_2^L-S(\xi_2,\xi_1) \right)} \\
        &=-\sum_{k,m=0}^\infty\oint_{C_{r_1}}d\xi_1\oint_{C_{R_2}}d\xi_2 \ \xi_1^{x_2-y_1-1+kL}\xi_2^{x_1-y_2-1-(m+1)L}S(\xi_2,\xi_1)^{k+m+2}.
    \end{split}
\end{equation}
 Using the change of variables $\xi_1=\eta/\xi_2$ we obtain
 \begin{equation}
     \begin{split}
         -\sum_{k,m=0}^\infty\oint_{C_{r=1/2}}d\eta\oint_{C_{R_2}}d\xi_2 \ \eta^{x_2-y_1-1+kL}\xi_2^{x_1-x_2+y_1-y_2-(k+m+1)(L-1)}\left(\frac{1/\xi_2+\eta/\xi_2-2\Delta}{1+\eta-2\Delta\eta/\xi_2}\right)^{k+m+2}.
     \end{split}
 \end{equation}
 Then we have the following upper bound on the $\xi_2$ exponent: $x_1-x_2+y_1-y_2-(k+m+1)(L-1)\leq -2-(k+m+1)(L-1)$. So there are no poles with respect to the $\xi_2$ variable and the above integrals are zero. Lastly we must compute the nested integral with $\xi_1$ over the $C_{R_1}$ contour and $\xi_2$ over the $C_{r_2}$ contour. We have 
 \begin{equation}
     \begin{split}
         &\oint_{C_{R_1}}d\xi_1\oint_{C_{r_2}}d\xi_2 \ \frac{\xi_1^{x_1-y_1-1}\xi_2^{x_2-y_2-1}}{\left( \xi_1^L-S(\xi_1,\xi_2) \right)\left( \xi_2^L-S(\xi_2,\xi_1) \right)} \\
         &= -\sum_{k,m=0}^{\infty}\oint_{C_{R_1}}d\xi_1\oint_{C_{r_2}}d\xi_2 \ \xi_1^{x_1-y_1-1-(k+1)L}\xi_2^{x_2-y_2-1+mL}S(\xi_1,\xi_2)^{k+m+1}
     \end{split}
 \end{equation}
 since $|\xi_1^{-L}S(\xi_1,\xi_2)|\approx \frac23|\Delta|\epsilon^{L-1}<1$ and $|\xi_2^LS(\xi_1,\xi_2)|\approx \frac43|\Delta| (2\epsilon)^{L-1}<1$. Using the change of variables $\xi_2=\eta/\xi_1$ we obtain 
 \begin{equation}
     \begin{split}
         -\sum_{k,m=0}^{\infty}\oint_{C_{R_1}}d\xi_1\oint_{C_{r=2}}d\eta \ \xi_1^{x_1-x_2+y_2-y_1-1-(k+m+1)(L-1)}\eta^{x_2-y_2-1+mL}\left(\frac{1/\xi_1+\eta/\xi_1-2\Delta}{1+\eta-2\Delta\eta/\xi_1}\right)^{k+m+1}.
     \end{split}
 \end{equation}
 On the exponent of $\xi_1$ we have the following upper bound: $x_1-x_2+y_2-y_1-1-(k+m+1)(L-1) \leq -2 +(L-1) -(k+m+1)(L-1)\leq -2$. Thus there are no poles with respect to $\xi_1$ so the above integrals are zero. On to the permuted terms, we have
 \begin{equation}
     \begin{split}
         &\oint_{C_{R_1}}d\xi_1\oint_{C_{r_2}}d\xi_2 \ \frac{S(\xi_2,\xi_1)\xi_1^{x_2-y_1-1}\xi_2^{x_1-y_2-1}}{\left( \xi_1^L-S(\xi_1,\xi_2) \right)\left( \xi_2^L-S(\xi_2,\xi_1) \right)} \\
         &= -\sum_{k,m=0}^{\infty}\oint_{C_{R_1}}d\xi_1\oint_{C_{r_2}}d\xi_2 \ \xi_1^{x_2-y_1-1-(k+1)L}\xi_2^{x_1-y_2-1+mL}S(\xi_1,\xi_2)^{k+m}.
     \end{split}
 \end{equation}
Using the change of variables $\xi_2=\eta/\xi_1$, we obtain
\begin{equation}
    \begin{split}
    \sum_{k,m=0}^{\infty}\oint_{C_{R_1}}d\xi_1\oint_{C_{2}}d\eta \ \xi_1^{x_2-x_1+y_2-y_1-1-(k+m+1)L - 1 + k + m}\eta^{x_1-y_2-1+mL}\left(\frac{1/\xi_1+\eta/\xi_1-2\Delta}{1+\eta-2\Delta \eta/\xi_1} \right)^{k+m}.
    \end{split}
\end{equation}
 We have the following upper bound on the $\xi_1$ exponent: 
 $x_2-x_1+y_2-y_1-(k+m)(L-1)-L-1
 \leq (2-k-m)(L-1)-L-1$.
 This upper bound is only positive for $k=m=0$. Otherwise, the exponent of $\xi_1$ is negative and the above integrals with respect to $\xi_1$ have no poles and are zero. For $k=m=0$, we are left with the nested integral 
 \begin{equation}
     -\oint_{C_{R_1}}d\xi_1\oint_{C_{r_2}}d\xi_2 \ \xi_1^{x_2-y_1-1-L}\xi_2^{x_1-y_2-1}.
 \end{equation}
 Then we deform the $C_{R_1}$ contour to $C_{r_1}$ and obtain 
\begin{equation}
    \oint_{C_{r_1}}d\xi_1\oint_{C_{r_2}}d\xi_2 \ \xi_1^{x_2-y_1-1-L}\xi_2^{x_1-y_2-1} = \mathds{1}(x_2-L=y_1, \ x_1=y_2).
\end{equation}

To piece everything together now, from evaluating our nested integrals with respect to poles inside the $C_1$ and $C_2$ contours, we have the following equalities that we've shown 
\begin{equation}
    \begin{split}
        \oint_{C_{r_1}}d\xi_1\oint_{C_{r_2}}d\xi_2 I(\boldsymbol{\xi};\mathbf{x},\mathbf{y}) &= \sum_{k=0}^\infty \mathds{1}(x+kL=y) + \sum_{k=1}^\infty\mathds{1}(x_2+(k-1)L=y_1,\ x_1+kL=y_2) \\
        \oint_{C_{R_1}}d\xi_1\oint_{C_{R_2}}d\xi_2 I(\boldsymbol{\xi};\mathbf{x},\mathbf{y}) &= \sum_{k\leq -1} \mathds{1}(x+kL=y) + \mathds{1}(x_2+(k-1)L=y_1,\ x_1+kL=y_2) \\
        \oint_{C_{r_1}}d\xi_1\oint_{C_{R_2}}d\xi_2 I(\boldsymbol{\xi};\mathbf{x},\mathbf{y}) &= 0 \\
        \oint_{C_{R_1}}d\xi_1\oint_{C_{r_2}}d\xi_2 I(\boldsymbol{\xi};\mathbf{x},\mathbf{y}) &= \mathds{1}(x_2-L=y_1,\ x_1=y_2).
    \end{split}
\end{equation}
Hence we have shown that 
\begin{equation}
\begin{split}
    \oint_{C_1}d\xi_1\oint_{C_2}d\xi_2 \ I(\boldsymbol{\xi};\mathbf{x},\mathbf{y}) &= \sum_{k\in\mathbb{Z}}\mathds{1}(x+kL=y) + \mathds{1}(x_2+(k-1)L=y_1,\ x_1+kL=y_2)\\
    &= \mathds{1}(x=_L y).
\end{split}
\end{equation}
\end{proof}

\end{document}
